% Editorial correspondence: gestion/mapa_transportes_auxiliares.md.
% Auxiliary realizations do not replace the main compatibility and reconstruction argument.
\section{Compensated transport, geometric memory and energy tests}
\label{md:aux:transporte}

\subsection{Operator origin and domain of the realization}
\label{md:aux:origen}

The construction retains APP--TRIT--TPK and the enriched state.
APP brings together the arithmetic sheets with their residues and
quotients; TRIT retains regime and orientation; TPK composes selection,
transport and memory updating. The matrix representations used below
are downstream realizations. The nonadic holonomy $\Gamma_9$ and its
memory retain their own operators.

The antecedents used are the elliptic, hyperbolic and parabolic
quadratic representations of TRIT, the differential of invertible
transports, the energy functional $E=\langle Df,WDf\rangle$ with
$W>0$, return with translational memory, the action ellipse with
semiaxis product $2\hbar$ and its LC realization, vacancy polarization
and reversible two-mode coupling. Translational return and the
commutator below are distinct objects. The elliptic--parabolic loop
and its oriented-area balance will subsequently be used as a test
collection.

The compensated composition and test protocol bring together these
realizations. The algebraic theory of commutator traces in
$\operatorname{SL}(2)$ is an established mathematical
antecedent.\footnote{Goldman, \emph{Trace Coordinates on
Fricke spaces}, \href{https://arxiv.org/abs/0901.1404}{arXiv:0901.1404}.}
The subject of this appendix is its specific composition with the
declared readout operators. The parameters $a,b$ are realization
parameters, not fitted physical constants.

The scope ends at this local realization. Identification of TPK routes
admitting a material device with those parameters requires composition
of their effective readout operators. An arbitrary invertible matrix
does not by itself establish a TPK route, and identification of the
commutation defect with Cartan torsion requires its connection and
soldering.

\subsection{Compensation of operations and area conservation}
\label{md:aux:conmutador}

The regime matrices are
\begin{equation}
 J_{\mathrm e}=\begin{pmatrix}0&1\\-1&0\end{pmatrix},\qquad
 J_{\mathrm h}=\begin{pmatrix}0&1\\1&0\end{pmatrix}.
 \label{md:aux:regimenes}
\end{equation}
They satisfy $J_{\mathrm e}^2=-I$, $J_{\mathrm h}^2=I$ and
$[J_{\mathrm e},J_{\mathrm h}]=2\operatorname{diag}(1,-1)$.
Define
\begin{equation}
 R(a)=e^{aJ_{\mathrm e}},\qquad B(b)=e^{bJ_{\mathrm h}},\qquad
 C_{\mathrm{aux}}(a,b)=R(a)B(b)R(a)^{-1}B(b)^{-1}.
 \label{md:aux:def-compensado}
\end{equation}
The action on columns is read from right to left. The protocol contains
each transformation and its inverse, but the sequence is not a full
retracing of the path. A full retracing would, for example, have the
form $RBB^{-1}R^{-1}=I$. The symbol $C_{\mathrm{aux}}$ distinguishes
this elliptic--hyperbolic commutator from the counterangle $C^*$ and
from the conjugate-shape loop $H_{\mathrm{lazo}}$ studied in the main
text.

\begin{proposition}
\label{md:aux:prop-traza}
For $a,b\in\mathbb R$,
\begin{equation}
 \det C_{\mathrm{aux}}=1,\qquad
 \operatorname{tr}C_{\mathrm{aux}}=2+4\sin^2a\sinh^2b.
 \label{md:aux:traza-compensado}
\end{equation}
\end{proposition}
\begin{proof}
Put $c_a=\cos a$, $s_a=\sin a$, $u_b=\cosh b$ and $v_b=\sinh b$.
Then
\[
 R=\begin{pmatrix}c_a&s_a\\-s_a&c_a\end{pmatrix},\qquad
 B=\begin{pmatrix}u_b&v_b\\v_b&u_b\end{pmatrix},
\]
with $c_a^2+s_a^2=1$ and $u_b^2-v_b^2=1$. All the matrices have
determinant one. Multiplication gives
\begin{equation}
 \begin{aligned}
 (C_{\mathrm{aux}})_{11}&=u_b^2+2c_as_au_bv_b-(c_a^2-s_a^2)v_b^2,\\
 (C_{\mathrm{aux}})_{12}&=-2s_a^2u_bv_b-2c_as_av_b^2,\\
 (C_{\mathrm{aux}})_{21}&=-2s_a^2u_bv_b+2c_as_av_b^2,\\
 (C_{\mathrm{aux}})_{22}&=u_b^2-2c_as_au_bv_b-(c_a^2-s_a^2)v_b^2.
 \end{aligned}
 \label{md:aux:entradas-compensado}
\end{equation}
The diagonal sum is
$2u_b^2-2(c_a^2-s_a^2)v_b^2=2+4s_a^2v_b^2$, as claimed.
\end{proof}

When $s_av_b\ne0$, the trace exceeds two and the commutator is
hyperbolic, with reciprocal positive eigenvalues
\begin{equation}
 e^{\lambda},\ e^{-\lambda},\qquad
 \lambda=2\operatorname{arsinh}|\sin a\sinh b|>0.
 \label{md:aux:espectro-compensado}
\end{equation}
This follows from $\cosh\lambda=\operatorname{tr}(C_{\mathrm{aux}})/2$
and $\cosh(2\operatorname{arsinh}x)=1+2x^2$.
If $s_av_b=0$, then $b=0$ or $R=\pm I$, and
$C_{\mathrm{aux}}=I$.
The local expansion is
\begin{equation}
 C_{\mathrm{aux}}=I+2ab\operatorname{diag}(1,-1)
 +O(|a|^2|b|+|a||b|^2).
 \label{md:aux:expansion-compensado}
\end{equation}
The residual begins with the product of the two operations. Applying
either operation followed by its own inverse yields the identity.

A two-dimensional real matrix of determinant one preserves the area
form. Thus $C_{\mathrm{aux}}$ transforms the action ellipse into
another ellipse of area $h$. If $\Sigma>0$ is a covariance matrix,
\begin{equation}
 \Sigma'=C_{\mathrm{aux}}\Sigma C_{\mathrm{aux}}^{\mathsf T},
 \qquad \det\Sigma'=\det\Sigma.
 \label{md:aux:covarianza-area}
\end{equation}
The distribution and orientation of the semiaxes may change.
The eigenvalues $e^{\pm\lambda}$ describe eigendirections and
iterations; singular values determine a different readout and do not,
in general, coincide with them. The variation of Euclidean semiaxes
in a single application requires the corresponding metric.

Conservation of action area is compatible with accumulation of
oriented deformation. The geometric observable $\lambda$ depends on
$a,b$ and does not constitute an additional universal constant.
Conservation of determinant or area alone does not establish
conservation of physical energy.

\subsection{Positive readout of the residual and path reversal}
\label{md:aux:residuo}

On a return fibre, with declared field $f$ and Hermitian weight $W>0$,
the energy readout operator takes the form
\begin{equation}
 E_C(f)=\langle(I-C_{\mathrm{aux}})f,
                  W(I-C_{\mathrm{aux}})f\rangle.
 \label{md:aux:energia-residuo}
\end{equation}
For $s_av_b\ne0$, the commutator has no eigenvalue one and
$E_C(f)>0$ for every $f\ne0$. Full retracing yields zero residual.
Dimensional normalization belongs to the initial functional: its
numerical coordinates require their realization before being
interpreted as joules, heat, mass or conduction current.

Complete path reversal inverts $C_{\mathrm{aux}}$.
The trace and $\lambda$ assign the same readout to the operator and its
inverse. Reconstruction of direction requires an additional readout:
the residual vector, transport components or a phase-resolved probe.
This comparison places a concrete limit on the sufficiency of an
exclusively scalar readout.

A rational example, independent of physical data, is
\begin{equation}
 R=\begin{pmatrix}3/5&4/5\\-4/5&3/5\end{pmatrix},\qquad
 B=\begin{pmatrix}5/4&3/4\\3/4&5/4\end{pmatrix}.
 \label{md:aux:ejemplo-racional}
\end{equation}
It yields $\operatorname{tr}C_{\mathrm{aux}}=86/25$,
$\det C_{\mathrm{aux}}=1$ and eigenvalues
$(43\pm6\sqrt{34})/25$. The rational checks in the supplement include
this example and $169$ pairs, covering zero cases, full retracing and
preservation of the covariance determinant. These are finite checks
of the proven identities.

\subsection{Reversible coupling and reduced entanglement readout}
\label{md:aux:acoplamiento}

For $H\in\operatorname{Sp}(2,\mathbb R)$, the two-mode realization
uses the coupling
\begin{equation}
 U_H=\frac19\begin{pmatrix}
 8I+H&\sqrt8(H-I)\\
 \sqrt8(H-I)&I+8H
 \end{pmatrix}.
 \label{md:aux:acoplamiento-uh}
\end{equation}
With two identical pure Gaussian states prepared from the action
ellipse and joint transport of the initial chart, the reduced readout
satisfies
\begin{equation}
 \nu(H)^2=1+\frac8{81}
    \bigl\{\operatorname{tr}(HH^{\mathsf T})-2\bigr\},
 \qquad\operatorname{Tr}\rho_{\mathrm{vis}}^2=\nu(H)^{-1}.
 \label{md:aux:mezcla-generica}
\end{equation}
Preparation and readout operators remain fixed. The joint
transformation is reversible and preserves global purity.

\begin{corollary}
\label{md:aux:cor-mezcla}
For $C_{\mathrm{aux}}(a,b)$,
\begin{equation}
 \operatorname{tr}(C_{\mathrm{aux}}C_{\mathrm{aux}}^{\mathsf T})-2
 =16\sin^2a\sinh^2b\cosh^2b,
 \label{md:aux:traza-cuadratica}
\end{equation}
and
\begin{equation}
 \boxed{\nu_C^2=1+\frac{32}{81}\sin^2a\sinh^2(2b).}
 \label{md:aux:mezcla-compensado}
\end{equation}
\end{corollary}
\begin{proof}
The four entries of \eqref{md:aux:entradas-compensado} can be
written as $d+k$, $l-m$, $l+m$, $d-k$, where
\[
 d=1+2s_a^2v_b^2,\quad k=2c_as_au_bv_b,\quad
 l=-2s_a^2u_bv_b,\quad m=2c_as_av_b^2.
\]
Their sum of squares is
$2(d^2+k^2+l^2+m^2)=2+16s_a^2u_b^2v_b^2$.
Substitution in \eqref{md:aux:mezcla-generica} and
$\sinh(2b)=2\sinh b\cosh b$ prove the result.
\end{proof}
Example \eqref{md:aux:ejemplo-racional} gives
$\nu_C^2=17/9$ and purity $3/\sqrt{17}$.

Cancellation of $\hbar$ and of the initial deformation uses the
declared preparation and joint chart transport.
A protocol with $s_av_b\ne0$ produces reduced mixedness; complete
inversion of the coupling can undo it. Reduced mixedness, stored
energy and heat retain their roles as distinct readouts.

\subsubsection{Signed test in an elliptic--parabolic collection}
\label{md:aux:familia-hn}

The constructed loop collection is
\begin{equation}
 H_n=\begin{pmatrix}1+n&n^2\\n&n^2-n+1\end{pmatrix},
 \qquad n\in\mathbb Z.
 \label{md:aux:hn}
\end{equation}
Its alternating memory balance depends on $n^2$, whereas the
covariance-determinant ratio is
\begin{equation}
 \mathcal R_n=\frac{\det V_{\mathrm{vis},n}}{\det V_0}
 =1+\frac8{81}n^2(2n^2-2n+5).
 \label{md:aux:covarianza-hn}
\end{equation}
Subtracting the two orientations cancels the even terms:
\begin{equation}
 \boxed{\mathcal R_{-n}-\mathcal R_n=\frac{32n^3}{81}.}
 \label{md:aux:diferencia-direccional}
\end{equation}
For $n=1$ and $n=-1$, the responses are $121/81$ and $153/81$,
with purities $9/11$ and $3/\sqrt{17}$, respectively.
This collection satisfies $H_{-n}\ne H_n^{-1}$ in general;
it is not the iteration $H_1^n$ either. Changing the oriented
parameter must be distinguished from fully retracing the operations.
Its alternating area readout retains an identification distinct
from that of the Barbero functional.

The mixedness criterion uses the commutator $[D,J]$, with
$D=\sqrt8(H-I)/9$. If $D$ commutes with the complex structure
compatible with the initial state, it may transport memory without
making that Gaussian product mixed. If it does not commute, the
readout operator determines the mixedness. Absence of mixedness
for that state is a property distinct from absence of material
dissipation.

The difference \eqref{md:aux:diferencia-direccional} is a corollary
of the collection and its readout operator, not a universal constant.
The formula for $C_{\mathrm{aux}}$ composes the same readout
operator with another collection admitted by the theorem. A material
test must determine transports and couplings before measurement;
arbitrarily programming these matrices establishes their emulation
and retains a different scope.

\subsection{Second moment of the vacancy current}
\label{md:aux:vacancias}

The capacity readout operator fixes
\[
 \delta=\log_{10}(10/9),\qquad z_n\in\{0,1\},\qquad
 P_N=\sum_{n<N}(z_n-\delta),\qquad |P_b-P_a|<1.
\]
On equal intervals $t_0$, the centred current is
$I_n=A_I(z_n-\delta)$, with $A_I=u_Q/t_0$.
The symbol $A_I$ distinguishes this amplitude from the principal
angle. For $L=b-a$,
\begin{equation}
 \frac1{A_I^2}\sum_{n=a}^{b-1}I_n^2
 =L\delta(1-\delta)+(1-2\delta)(P_b-P_a).
 \label{md:aux:segundo-momento}
\end{equation}
The proof expands $(z_n-\delta)^2$, uses $z_n^2=z_n$ and sums.
The signed imbalance remains bounded, whereas
$I_{\mathrm{rms}}/|A_I|\longrightarrow\sqrt{\delta(1-\delta)}$.
The amplitude $A_I$ is the jump between the two current levels.
The test retains the temporal realization and must account for
any coupling that filters the sequence.

A zero signed balance can coexist with a positive quadratic readout.
If an ideal resistance $R_E>0$ is added and that current is
maintained over the intervals, the downstream test model gives
\begin{equation}
 Q_{R_E}=R_Et_0\sum_n I_n^2.
 \label{md:aux:contraste-resistivo}
\end{equation}
The resistance belongs to this downstream realization. Its role is
to test a possible heat reduction against an explicit balance.

\subsection{Conditions for a material test}
\label{md:aux:contraste-material}

The protocol separates two questions. For order memory, compensated
paths, order exchanges and full retracings are compared through a
remaining vector or phase-resolved response. An isolated intensity
or net charge may lose the distinction sought. For energy recovery,
input, useful or returned energy, storage variation and heat are
measured, including preparation and restoration of the apparatus.
Phase return remains distinct from variation in stored energy.

In a scalar one-port test, the relation
$v_2(t)=v_1(T-t)$ preserves the amplitude distribution and power
spectrum. In a periodic regime, the same linear time-invariant system
predicts equal mean power. An instantaneous law $i=g(v)$ also preserves
the integral of $vi$. A reproducible difference would probe history
dependence, but attributing it requires comparison with hysteresis,
heating and alternative dynamical models. In multiport systems,
cross-spectra must also be retained.

A reduction in dissipation is tested at equal useful transfer,
fidelity and rate. Attenuation must distinguish energy, signal
amplitude, carrier population and conserved charge before it is
attributed to a gravitational origin. The distinction between storage,
recovery and losses has experimental precedents in electrocaloric
recovery.\footnote{Defay et al., \emph{Enhanced electrocaloric
efficiency via energy recovery} (2018),
\href{https://doi.org/10.1038/s41467-018-04027-9}{doi:10.1038/s41467-018-04027-9}.
This antecedent does not constitute experimental evidence for the
HMT transport composed here.}

The appendix yields a general identity within the specified matrix
collection, a positive readout operator for the residual, an inversion
check and an exact law for the second moment of the current.
Its scope distinguishes these results from a new universal constant,
a prototype without heating or a proof of gravitational electron
decay. The prospective link with a device must be fixed before
obtaining the test data.

\section{Action, reduced memory and recoverable work}
\label{md:aux:trabajo}

\subsection{Composition of action with an energy realization}
\label{md:aux:accion}

The composition receives action outputs and memory transports from
APP--TRIT--TPK, retaining the enriched state as an antecedent.
Its domain is the declared Gaussian and energy realization.
Passivity and ergotropy are established concepts; they are used to
study the specific coupling and the readout operators already
constructed.

The action equation keeps its coefficients explicit:
\begin{equation}
 H_5=\frac12\sqrt{\frac{1000\alpha}{\varphi}}-\alpha
 +\frac9{16}\alpha^2-\frac59\alpha^3
 +\frac7{48}\alpha^4-\frac1{54}\alpha^5,
 \label{md:aux:h5}
\end{equation}
\begin{equation}
 D_A=e^{-100\pi\alpha/9},\qquad
 \mathcal C_\pi=169\alpha^6+\frac{D_A\alpha^7}{1-D_A\alpha},
 \label{md:aux:correccion-accion}
\end{equation}
\begin{equation}
 \begin{aligned}
 \hbar_{\mathrm{pre}}&=s_0H_5,\\
 \hbar_{\mathrm{ret}}&=s_0\left(H_5-\frac{90}{\pi}\mathcal C_\pi\right),
 \qquad h_s=2\pi\hbar_s,\qquad s_0=10^{-34}\mathcal U_S.
 \end{aligned}
 \label{md:aux:secciones-accion}
\end{equation}
The decimal order and dimensional chart retain distinct antecedents.
The formula keeps the coordinates $\pi,\varphi$, the orientation
and scale of the readout operator, in addition to $\alpha$.
The dimensional equality retains its unit of action.

Fix a positive section and write $\hbar=\hbar_s$.
The ellipse has semiaxes $\sqrt{2\hbar}e^{\zeta/2}$ and
$\sqrt{2\hbar}e^{-\zeta/2}$, where
$\zeta=\operatorname{artanh}(C^*/A)$. Its oriented area satisfies
\begin{equation}
 \mathcal A(\theta)=\hbar\theta,\qquad
 \mathcal A(2\pi)=h.
 \label{md:aux:area-fase}
\end{equation}
In the LC realization, with declared clock $\omega$ and reference
impedance,
\begin{equation}
 \mathcal H_\zeta(Q,P)
 =\frac\omega2\left(e^{-\zeta}Q^2+e^\zeta P^2\right).
 \label{md:aux:hamiltoniano-lc}
\end{equation}
On the orbit, $\mathcal H_\zeta=\hbar\omega=hf$, with
$f=\omega/(2\pi)$. The Hamiltonian convention gives
$\dot\theta=-\omega$, and hence
$d\mathcal A/dt=-\hbar\omega$. Positive energy coincides with the
magnitude of that areal velocity in the realization and fixed section
considered. Accumulated area and the Lagrangian action of an arbitrary
trajectory are distinct objects.

The internal equation determines the action scale per unit phase;
the clock fixes its rate. Ellipse deformation and transport order
retain information beyond $hf$. The Gaussian preparation used
subsequently differs from a point on the classical orbit: its initial
energy is $hf/2$, rather than $hf$.

\subsection{Preparation and readout operators for the energy result}
\label{md:aux:hipotesis-gaussianas}

Two pure, centred, identical Gaussian modes are prepared with
covariance
\begin{equation}
 V_0=\frac\hbar2MM^{\mathsf T},\qquad
 M=\operatorname{diag}(e^{\zeta/2},e^{-\zeta/2}).
 \label{md:aux:preparacion}
\end{equation}
The operator $U_H$ of \eqref{md:aux:acoplamiento-uh} acts,
transported to the initial chart by $M\oplus M$. The joint state
remains pure. The final energy readout operators have the same
frequency $\omega$, retain the initial metric
$G_\zeta=M^{-\mathsf T}M^{-1}$ and include no residual interaction
energy between the modes. Extraction operations begin and end with
these same readout operators. Physical accessibility of the admitted
operations is a hypothesis of the work bound, distinct from
construction of a device realizing them.

Define
\begin{equation}
 \begin{gathered}
 \tau=\operatorname{tr}(HH^{\mathsf T})-2\geq0,\\
 W_v=\frac{8I+HH^{\mathsf T}}9,\qquad
 W_m=\frac{I+8HH^{\mathsf T}}9.
 \end{gathered}
 \label{md:aux:covarianzas-reducidas}
\end{equation}
The covariances are $V_j=(\hbar/2)MW_jM^{\mathsf T}$.
Visible purity is denoted by $\mu=\operatorname{Tr}\rho_v^2$
and its normalized symplectic eigenvalue by $\nu=\mu^{-1}$.
In this appendix, $\nu$ has this role; ordinary frequency is written
$f$ and angular frequency $\omega$.

\subsection{Identity between visible energy and purity}
\label{md:aux:energia-pureza}

\begin{proposition}
\label{md:aux:prop-energia-pureza}
With the preceding preparation and readout operators, the visible
energy above the initial state satisfies
\begin{equation}
 \frac{\Delta E_v}{hf}=\frac\tau{36},\qquad
 \nu^2=1+\frac{8\tau}{81},\qquad
 \boxed{\mu^{-2}-1=\frac{32}{9}\frac{\Delta E_v}{hf}.}
 \label{md:aux:identidad-energia-pureza}
\end{equation}
\end{proposition}
\begin{proof}
The state is centred and the quadratic readout operator remains fixed.
Hence,
\[
 E_v=\frac\omega2\operatorname{tr}(G_\zeta V_v)
 =\frac{\hbar\omega}{4}\operatorname{tr}W_v.
\]
Since $\operatorname{tr}W_v=2+\tau/9$, subtracting
$\hbar\omega/2$ gives the first equality.
The reduced covariance satisfies
$\det W_v=1+8\tau/81$ and $\mu=(\det W_v)^{-1/2}$.
Eliminating $\tau$ proves the boxed identity.
The proof requires neither symmetry of $H$ nor coincidence between
its eigenvalues and singular values.
\end{proof}

The law eliminates the particular path parameters.
The coefficient $32/9$ comes from the coupling split, and the scale
$h$ is retained until normalization.
To check the specificity of that split, consider
$W_p=pI+(1-p)HH^{\mathsf T}$ with $0<p<1$. One obtains
\begin{equation}
 \mu^{-2}-1=4p\frac{\Delta E}{hf}.
 \label{md:aux:reparto-general}
\end{equation}
The declared choice is $p=8/9$. When $\Delta E>0$, a test can estimate
\begin{equation}
 p_{\mathrm{obs}}=
 \frac{\mu^{-2}-1}{4\Delta E/(hf)}
 \label{md:aux:reparto-observado}
\end{equation}
and compare it with $8/9$ without fitting $p$ to those same data.
Programming a mixer with that split verifies an emulation;
attribution of the split to a natural mechanism retains its
independent proof.

The identity uses centred Gaussian states and the specified
preparation. A coherent displacement adds energy without changing
purity. Initial thermal noise and metric changes alter the relation.
These variations delimit the statement and provide falsifiers for
an extension to arbitrary states.

The convex combination $W_v$ combines covariance matrices produced
by reducing a joint Gaussian state. A probabilistic mixture of two
pure states is a different object, generally non-Gaussian; its purity
is not given by the determinant of that covariance.

\subsection{Recoverable work under local operations}
\label{md:aux:ergotropia}

Ergotropy is the maximum energy decrease under cyclic unitary
operations with a fixed energy readout operator.
Its definition does not automatically include the net cost of
constructing, preparing or controlling the apparatus.

The reduced spectrum has populations
\begin{equation}
 \lambda_j=(1-r)r^j,\qquad r=\frac{\nu-1}{\nu+1}.
 \label{md:aux:espectro-gaussiano}
\end{equation}
They decrease with oscillator energy.
The passive state therefore has energy $E_{v,\mathrm{pas}}=hf\nu/2$.
It is reached by a Gaussian operation:
\begin{equation}
 S=\sqrt\nu\,W_v^{-1/2},\qquad
 \det S=1,\qquad SW_vS^{\mathsf T}=\nu I.
 \label{md:aux:pasivacion}
\end{equation}
Population ordering also proves minimality among all unitaries,
not only Gaussian ones.

Consequently,
\begin{equation}
 \boxed{
 \frac{\mathcal W_v}{hf}
 =\frac{(\nu-1)(9\nu-7)}{32},\qquad
 \frac{E_{v,\mathrm{pas}}-hf/2}{hf}=\frac{\nu-1}{2}.}
 \label{md:aux:trabajo-visible}
\end{equation}
The two quantities sum to
$\Delta E_v/(hf)=9(\nu^2-1)/32$. The passive term is energy
inaccessible under the declared local operations; its presence does
not amount to an actual heat flow.

\subsubsection{Joint access to correlations}
\label{md:aux:acceso-conjunto}

The second mode satisfies $\det W_m=\det W_v=\nu^2$ and
$\Delta E_m=8\Delta E_v$. Thus,
\begin{equation}
 \Delta E_{\mathrm{tot}}=9\Delta E_v=hf\frac\tau4.
 \label{md:aux:energia-conjunta}
\end{equation}
Joint inversion of the coupling restores the initial product.
The total state is pure, and that product is the ground state of
the joint readout operator; joint access ideally allows recovery of
$\Delta E_{\mathrm{tot}}$.
Independent local operations leave two passive states, each with
energy $hf\nu/2$. It follows that
\begin{equation}
 \boxed{\mathcal W_{\mathrm{conjunta}}-\mathcal W_{\mathrm{local}}
 =hf(\nu-1)=hf(\mu^{-1}-1).}
 \label{md:aux:ventaja-conjunta}
\end{equation}

The identity quantifies the additional work accessible through
correlations. The initial preparation required energy, and the
balance of a complete cycle includes that contribution.
The relation compares energy readouts and operations on the state,
without identifying information and energy as the same physical
quantity.

Here “local” means product unitaries $U_v\otimes U_m$;
measurements, classical communication and feedback lie outside that
class. The gain follows from the operational restriction on the
correlated state and from its marginal spectra, not from residual
interaction energy. The $1:8$ split concerns energies above the vacuum,
not total energies or ergotropies.

\subsection{Exact cases and transport orientation}
\label{md:aux:casos-energia}

For $H_n$ in \eqref{md:aux:hn},
\begin{equation}
 \tau_n=n^2(2n^2-2n+5).
 \label{md:aux:tau-hn}
\end{equation}
At $n=1$, one obtains
\begin{equation}
 \begin{gathered}
 \nu=11/9,\qquad\mu=9/11,\qquad
 \Delta E_v=\frac5{36}hf,\\
 \mathcal W_v=\frac1{36}hf,\qquad
 E_{v,\mathrm{pas}}-hf/2=\frac19hf.
 \end{gathered}
 \label{md:aux:caso-positivo}
\end{equation}
The joint system has $\Delta E_{\mathrm{tot}}=5hf/4$, local work
$37hf/36$ and joint advantage $2hf/9$.
The case $n=-1$ shares the eigenvalues of $H_1$ but has $\tau=9$ and
\begin{equation}
 \nu=\frac{\sqrt{17}}3,\qquad
 \Delta E_v=\frac{hf}{4},\qquad
 \mathcal W_v=\frac{9-2\sqrt{17}}{12}hf.
 \label{md:aux:caso-negativo}
\end{equation}
Added visible energy differs by a ratio $9/5$ at equal frequency
and action section. The change $n\mapsto-n$ is not full retracing:
$H_{-n}\ne H_n^{-1}$ in general.
True matrix inversion preserves $\tau$ in dimension two.
The parametric asymmetry of this collection thus remains distinct from
complete reversal.

For $C_{\mathrm{aux}}$,
$\tau_C=4\sin^2a\sinh^2(2b)$. The energy composition gives
\begin{equation}
 \boxed{\Delta E_v(C_{\mathrm{aux}})
 =\frac{hf}{9}\sin^2a\sinh^2(2b),}
 \label{md:aux:energia-compensado}
\end{equation}
\begin{equation}
 \mathcal W_{\mathrm{conjunta}}-\mathcal W_{\mathrm{local}}
 =hf\left(\sqrt{1+\frac{32}{81}\sin^2a\sinh^2(2b)}-1\right).
 \label{md:aux:trabajo-compensado}
\end{equation}
The connection with Planck remains explicit: in the prior section,
$hf=2\pi s_0H_5(\alpha,\varphi)f$; in the return section,
$H_5$ is replaced by $\eta_{\mathrm{ret}}$.
The composition fixes the section used; it assumes neither an
experimental variation of $\alpha$ nor free alternation between
dimensional sections.

\subsection{Boltzmann and the equivalent temperature of the passive state}
\label{md:aux:temperatura}

For $\nu>1$, the passive state has equivalent temperature
\begin{equation}
 k_BT_{\mathrm{pas}}
 =\frac{hf}{\log[(\nu+1)/(\nu-1)]}.
 \label{md:aux:temperatura-pasiva}
\end{equation}
When the transducer clock is used,
$hf=k_B\Theta_{\mathrm{clk}}\log3$, giving
\begin{equation}
 \frac{T_{\mathrm{pas}}}{\Theta_{\mathrm{clk}}}
 =\frac{\log3}{\log[(\nu+1)/(\nu-1)]}.
 \label{md:aux:razon-temperaturas}
\end{equation}
For $n=1$, this recovers $\log3/\log10$.
The entropy is that of the Gaussian spectrum, with occupation
$(\nu-1)/2$.
The energy composition retains this previously constructed thermal
readout. The equivalent temperature of the passive state does not imply
that the preceding reduced state was in equilibrium or that heat has
been transferred.

\subsection{Specificity and testing conditions}
\label{md:aux:contraste-trabajo}

The general relation between coherence, correlations and work
extraction belongs to the existing literature.
The composition examined here specifies a joint law: coefficient
$32/9$, energy split $1:8$, oriented collection $H_n$, action scale and
the difference between local and joint access.

Two experimental levels are distinguished. Emulation realizes $U_H$
and checks its identities; it verifies that realization of the model.
Physical prediction requires prospective determination, from the
construction, of the device, degrees of freedom, frequency, preparation
and coupling realizing its objects, without fitting $8/9$ to the test
data. Energy, covariance, purity, control costs and restoration energy
are measured separately.

An independent test obtains temperature and purity through readout
operators independent of the law under examination. Recovered work is
compared with a balance that includes preparation, controls and
memory restoration.

External comparisons include studies of Gaussian ergotropy and
passivation,\footnote{Kua, Serafini and Genoni, \emph{Daemonic ergotropy
of Gaussian quantum states and the role of measurement-induced
purification via general-dyne detection},
\href{https://arxiv.org/abs/2506.22288}{arXiv:2506.22288}.
Its role is comparison with general formulas for energy, purity and
passivity; it does not validate the HMT coupling or its constants.}
coherent ergotropy extraction with explicit energy
costs,\footnote{\emph{Experimental Extraction of Coherent Ergotropy
and Its Energetic Cost in a Superconducting Qubit},
\href{https://arxiv.org/abs/2506.16881}{arXiv:2506.16881}.}
and coherent control between quantum-machine
cycles.\footnote{Hou et al., \emph{Combining energy efficiency and
quantum advantage in cyclic machines},
\href{https://doi.org/10.1038/s41467-025-60179-5}{doi:10.1038/s41467-025-60179-5}.
The cited experiment does not measure the nonadic transport composed
in this appendix.}

The result is an algebraic composition with explicit initial
structure. Its proofs establish identities and operational
restrictions. An observed experimental advantage, a new universal
constant, a violation of thermodynamics or a material mechanism for
heat-free transport would require their respective demonstrations
and tests; they are not conclusions of the auxiliary identities
established here.
